\section{Representation acquisition is upstream of the licence}\label{sec:acquisition-licence}

A transfer decision has two logically distinct stages. The first constructs candidate structure from the source and target material. The second decides whether a proposed mapping is licensed. Writing these stages separately prevents a capable verifier from being credited for a representation it was given, and prevents a failed extractor from being interpreted as a theorem that no useful mapping exists:
\[
(X_s,X_t)\xrightarrow{A}(R_s,R_t)\xrightarrow{M}W\xrightarrow{L}V.
\]
Here $A$ is a representation-acquisition operator, $M$ constructs a candidate mapping or morphism, $W$ is the evidence-bearing witness, and $L$ is the fail-closed licence studied in this paper. The verdict $V$ lies in \{\texttt{LICENSED},\texttt{REJECTED},\texttt{CANNOT\_CHECK}\} before the stronger impossibility subtype is considered.

The factorization is executor-independent. $A$ may be a human annotation procedure, a rule-based parser, graph induction, anti-unification, program synthesis, a theorem prover, an LLM, or a hybrid. Likewise, $M$ may be a classical graph matcher, a domain-specific mathematical map, or a learned proposal generator. The contract does not grant extra authority to a witness because an LLM produced it, and it does not discount a witness because a deterministic program produced it. What matters is whether the registered obligations are discharged.

This separation is consistent with long-standing work on analogy. Structure-mapping theory distinguishes relational correspondence from surface similarity \cite{gentner1983structure}, and the Structure-Mapping Engine turns that distinction into an explicit algorithmic mapping procedure \cite{falkenhainer1989sme}. Experiments on schema induction further show that comparison of multiple analogs can support transfer by inducing a more abstract schema \cite{gick1983schema}. These results motivate acquisition mechanisms that construct or revise a representation rather than treating structure as a field already present in the input.

They also explain why the natural-domain negative in this paper should be localized to the registered reducer. The ARN experiment used third-party labels and passed the instrument battery, but its admitted deterministic reducer did not recover the structural relations needed for cross-domain transfer. The correct conclusion is therefore an acquisition-capability negative for that reducer, not a negative theorem about the licence and not an impossibility theorem about analogical abstraction.

\subsection{Progressive abstraction and explanation as prospective acquisition mechanics}

The open Pursuit-Plane programme now treats representation construction as its own research object. One candidate is progressive relational abstraction: align highly comparable examples first, induce a candidate schema, test the schema on held-out or near-miss cases, and only then widen the semantic distance. This proposal has strong cognitive parents. Gick and Holyoak found that schemas induced from multiple analogs were predictive of later transfer \cite{gick1983schema}. Explanation can also bias learning toward broad, causally useful regularities rather than surface detail \cite{williams2010explanation,walker2014explanation,legare2014explanation}. These findings do not establish the RAKL candidate. They define stronger parents and falsifiers for it.

A second candidate is a typed morphism registry. A generic witness should identify the mathematical family of map it proposes rather than asking one informal relation to carry every meaning. Causal abstraction provides a concrete parent: exact and stronger transformations preserve specified interventional behavior across levels of description \cite{rubenstein2017causal,beckers2019abstracting}. A future registry may similarly dispatch causal, algebraic, refinement, dimensional, proof-translation or constraint-preserving witnesses to family-specific validators while the outer Paper-II contract supplies shared evidence, boundary, refusal and forbidden-loss semantics. This remains prospective until domain validators are implemented and tested.

The current reference runtime also contains a narrower formal-parent arbitration boundary. When a strongest formal parent returns a decisive in-scope verdict, that verdict dominates the generic Orion route; a formal 	exttt{CANNOT\_CHECK} cannot be overridden; Orion delegation is permitted only after an evidence-bound 	exttt{OUT\_OF\_SCOPE} result. This is an implementation and fail-closed-composition result, not evidence that the generic contract outperforms the formal parent. It is the pattern a future typed morphism registry should preserve.

\subsection{LLMs as acquisition adapters, not transfer authority}

LLMs are attractive at the acquisition stage because scientific structure is often described in prose rather than supplied as a typed graph. Recent tool-using scientific agents demonstrate that language models can coordinate retrieval, planning, code and experiments \cite{boiko2023coscientist,bran2024chemcrow,gottweis2026coscientist}. But the logical role is narrow: an LLM is an adapter that proposes $R_s$, $R_t$ or $W$. It is not the licence $L$ and it does not certify its own mapping.

This gives Paper~II an explicit LLM-free baseline programme. If source and target structures are already formal, the full licensing path can be deterministic. If they are not, candidate acquisition can be compared across human, symbolic, LLM and hybrid executors under the same downstream contract. The resulting experiment asks a sharper question than ``does the LLM reason analogically?'': which executor recovers a representation that survives the same typed obligations, at what cost, and on which scope?

\subsection{Selection processes and absent evidence}

Finally, acquisition is affected by how examples enter the corpus. Pedagogical-learning models show that the same examples can support different inferences when they were deliberately selected by a knowledgeable teacher rather than sampled independently \cite{shafto2014pedagogy}. Scientific corpora are also selected through publication, retrieval, citation and curation processes. Paper~II therefore treats future evidence selection as part of provenance. Absence of a counterexample or missing relation is informative only to the extent that the registered acquisition/selection process made it observable. Unknown selection remains a reason for \texttt{CANNOT\_CHECK}, not a free licence.
